\thispagestyle{empty}

\section*{Abstract}\label{sec:abstract}
This thesis documents the development and testing of the RaspberryPI Cyber Security Laboratory (RCSL).
The RCSL is imagined as a testing platform for demonstration and practice purposes of Cybersecurity in various IT applications.

An expandible software architecture is developed, which focuses on usability and includes features for the creation of test environments in the WiFi and web application domain.
Development and functionality are documented, highlighting the challenges encountered along the process.

Testing is conducted in the form of cracking the WEP and WPA/WPA2 network standards and the results are evaluated.
Along testing several faults in the system were discovered and mostly corrected.
The results show the RCSL to be capable of establishing testing environments for WiFi security.

Proposed usecases for the RCSL include demonstration and hands-on practice in cybersecurity and software development education.
Future improvement and expansion of the project are recommended.


\newpage

\section*{Kurzdarstellung}
\label{sec:kurzdarstellung}
Diese Arbeit dokumentiert die Entwicklung und Erprobung des RaspberryPI Cyber Security Laboratory (RCSL).
Das RCSL ist als Testplattform für Demonstrations- und Übungszwecke von Cybersecurity in verschiedenen IT-Anwendungen gedacht.

Es wird eine erweiterbare Software-Architektur entwickelt, die sich auf die Benutzerfreundlichkeit konzentriert und Funktionen für die Erstellung von Testumgebungen im WiFi- und Web-Anwendungsbereich enthält.
Entwicklung und Funktionalität werden dokumentiert und die dabei aufgetretenen Herausforderungen aufgezeigt.

Als Tests werden das cracking der Netzwerkstandards WEP und WPA/WPA2 durchgeführt und die Ergebnisse ausgewertet.
Während der Tests wurden mehrere Fehler im System entdeckt und größtenteils behoben.
Die Ergebnisse zeigen, dass das RCSL in der Lage ist, Testumgebungen für WiFi-Sicherheit einzurichten.

Vorgeschlagene Anwendungsfälle für den RCSL sind Demonstrationen und praktische Übungen in der Cybersecurity- und Softwareentwicklungslehre.
Zukünftige Verbesserungen und Erweiterungen des Projekts werden empfohlen.